% =============================================================
% Section 2. Related Work
% Template: paper/v0.1/notes/narrative.md §4 §2
% Structure: §2.1 FEEC / §2.2 cellular sheaf SP / §2.3 TopSP / §2.4 dagger Galerkin
% Key message: each school is precise; the four have not been combined.
% =============================================================

\noindent
The technical work in this paper combines tools from four schools
that have developed independently. Each school is internally precise;
none has been combined with the others on a single concrete object
producing a single scalar. We review the relevant results here so
that later sections can cite them by tag.

% -------------------------------------------------------------
\subsection{Finite element exterior calculus}\label{subsec:related-feec}

Finite element exterior calculus rests on three lines of
construction. The reduction--reconstruction interface of
\cite{BochevHyman2006} (with antecedents in
\cite{Whitney1957,Bossavit1988}) defines a map
$R : \Lambda^k(\Omega) \to C^k$ that integrates a smooth $k$-form
over chains, and an inverse $I_W : C^k \to \Lambda^k$ that lifts
chains to forms via Whitney shape functions; these two maps satisfy a
commuting Stokes identity (\cref{prop:commuting-stokes}). The
finite-element-system framework of \cite{Christiansen2010} packages
these data abstractly: an FES is a contravariant functor on a
category of cellular complexes, and Whitney forms are recovered as
its canonical instance. The Hilbert complex framework of
\cite{AFW2010,FalkWinther2014} (refining
\cite{BruningLesch1992,AFW2006}) supplies the analytic core that we
need: closed cochain complexes of Hilbert spaces, Hodge decomposition
with stalk-wise harmonic representatives, and bounded cochain
projections $\pi_h$ commuting with the differential up to a fixed
constant. The recent survey of \cite{ChristiansenRapetti2025}
collates 30 years of FEEC development and is one of our two anchors
for placing FEEC tools relative to the topological-signal-processing
literature.

% -------------------------------------------------------------
\subsection{Cellular sheaves and sheaf signal processing}\label{subsec:related-sheafsp}

Cellular sheaves of vector spaces, with cohomology computed by a
chain complex on a regular cell complex, are surveyed in
\cite[\S 6]{Robinson2014}. Robinson defines sheaf bandwidth as a
real-valued analytic invariant and closes the chapter with the open
question on its categorically intrinsic formulation that motivates
the present paper. An earlier sheaf-theoretic perspective on
sampling, complementary in topic to the bandwidth question studied
here, appears in \cite{Robinson2014Sampling}; the abstract sampling
cell complex $X^A$ and the strict-equality construction of
\cref{thm:strict-equality} are independent of that work, supplying a
specific finite-dimensional sampling sheaf on which Robinson's open
question admits a concrete spectral resolution. The spectral theory of cellular sheaf Laplacians
$L^k = (\delta^k)^* \delta^k + \delta^{k-1} (\delta^{k-1})^*$ is
developed in \cite{HansenGhrist2019}; we use Definition~1 of that
paper, especially Cond.~4 (boundaries of $1$-cells consist of two
distinct $0$-cells), to certify regularity of the abstract sampling
cell complex of \cref{def:Xsamp}. The conference
version \cite{HansenGhrist2019Allerton} introduces the spectral gap
quantities $\gamma_+, \Gamma_+$ on which our intrinsic invariant
$\sheafBW := \sqrt{\gamma_+/\Gamma_+}$ is built.

% -------------------------------------------------------------
\subsection{Topological signal processing}\label{subsec:related-topsp}

Topological signal processing on simplicial complexes is the program
initiated by \cite{BarbarossaSardellitti2020}: signals live on
$k$-cochains, the simplicial Laplacian replaces the graph Laplacian,
and bandlimited subspaces are defined as eigenspaces of the
simplicial Laplacian below a fixed threshold $B$. The school treats
sampling and reconstruction via Riesz frame inequalities in the
bandlimited subspace, an angle that we recover as condition~(b) of
\cref{thm:triangle-equiv}. The recent survey of
\cite{CerejeirasKahlerLegatiuk2025} is our second anchor for
placing topological signal processing relative to FEEC; that survey
does not invoke AFW Hilbert complexes.

% -------------------------------------------------------------
\subsection{Dagger Galerkin theory}\label{subsec:related-dagger}

A categorical formulation of Galerkin theory in which monomorphisms
carry a $\dagger$ (right-adjoint) structure is given by
\cite{LahtinenStenvall2020}; the abstract dagger setting is
\cite{HeunenVicary2019}. We use this thread sparingly: the dagger
tools clarify which subcategories of $\CellSh^{\fin}_{\Hilb}(X)$
carry a unitary structure, and they make
\cref{op:non-trivial-dagger} a precise question about the existence
of a non-trivial $\dagger$ on the full category. The unitary
subcategory $\CellSh^{\fin}_{\Hilb}(X)^{\mathrm{u}}$ of
\cref{def:unitary-subcat} is a groupoid
(\cref{prop:unitary-groupoid}); the question is whether this
restricted structure extends.

% -------------------------------------------------------------
\subsection{What has not been combined}\label{subsec:related-gap}

Each of the four schools above provides precise tools, and each
admits the underlying combinatorial object---a regular cell complex
with cochain spaces and a coboundary operator---without difficulty.
Yet the four have not been combined on a single concrete object
producing a single scalar that is at once intrinsic to the abstract
structure and computable from acquisition data. The references
cited in this section are mutually disjoint: \cite{AFW2010} does
not invoke cellular sheaves; \cite{HansenGhrist2019} does not
invoke FEEC; \cite{BarbarossaSardellitti2020} does not invoke
either; \cite{LahtinenStenvall2020} does not invoke any signal
processing applications. This paper supplies the missing object
(the abstract sampling cell complex $X^A$ of \cref{def:Xsamp}) and
the missing scalar (the strict equality of
\cref{thm:strict-equality}). \Cref{tab:feec-four-schools} in
\cref{sec:feec_review} records the four-school comparison once the
needed FEEC machinery has been reviewed.

\medskip\noindent
Each school provides precise tools but operates in isolation. We now
exhibit the first formal bridge between three of them via a triangle
equivalence theorem (\cref{sec:triangle}), which pinpoints exactly
where Robinson's open problem lives.
